% PAGES: 244-244

Note that
\begin{equation}
y^tAx \;=\; x^tA^ty
\end{equation}
since $(u,v) = (v,u)$ for any $u$ and $v$ and therefore
\[
y^tAx \;=\; (Ax,y) \;=\; (y,Ax) \;=\; (Ax)^ty \;=\; x^tA^ty.
\]

From (8.4) and (8.5), it follows that
\begin{equation}
||y-Ax||^2 \;=\; y^ty \;-\; 2x^tA^ty \;+\; x^tA^tAx.
\end{equation}

From (8.3) and (8.6), and since $y^ty = ||y||^2$, we conclude that
\begin{equation}
f(x) \;=\; ||y||^2 \;-\; 2x^tA^ty \;+\; x^tA^tAx.
\end{equation}

Recall that any minimum point $x_0$ of $f(x)$ must be a critical point of $f(x)$, i.e.,
a solution to $Df(x_0) = 0$. From (8.7), it follows that the gradient $Df(x)$ of $f(x)$
is
\begin{equation}
Df(x) \;=\; -2D(x^tA^ty) \;+\; D(x^tA^tAx),
\end{equation}
since $||y||^2$ is not a function of $x$ and therefore $D(||y||^2) = 0$.

Recall from (10.44) and (10.47) the following gradient formulas:
\begin{align}
D(x^tC) &= C^t \\
D(x^tMx) &= 2(Mx)^t,
\end{align}
for any constant column vector $C$ and for any symmetric matrix $M$. Since $A^tA$ is a
symmetric matrix, see Lemma 5.2, we obtain from (8.9) and (8.10) that
\begin{align}
D(x^tA^ty) &= (A^ty)^t; \\
D(x^tA^tAx) &= 2(A^tAx)^t.
\end{align}

From (8.8), (8.11), and (8.12), it follows that
\begin{align}
Df(x) \;&=\; -2(A^ty)^t + 2(A^tAx)^t \notag\\
&=\; 2(A^tAx - A^ty)^t.
\end{align}

From (8.13), we find that $Df(x_0) = 0$ if and only if
\begin{equation}
A^tAx_0 \;=\; A^ty.
\end{equation}

Since the columns of $A$ are linearly independent, it follows from Lemma 5.2 that the
matrix $A^tA$ is symmetric positive definite. Then, $A^tA$ is a nonsingular matrix, see
Lemma 5.3, and therefore the unique solution of (8.14) is
\begin{equation}
x_0 \;=\; (A^tA)^{-1}A^ty.
\end{equation}

To classify the critical point $x_0$ given by (8.15), we compute the Hessian of
$f(x)$. From (10.46) and (10.48), we obtain that $D^2(x^tA^ty) = 0$ and
$D^2(x^tA^tAx) = 2(A^tA)$. Then, from (8.7), it follows that $D^2f(x) = 2(A^tA)$, and
therefore
\[
D^2f(x_0) \;=\; 2(A^tA).
\]
